\chapter{Use Cases}
\ChapterLabel{UseCases}
\section{How to Use This Chapter}
Describe goal-oriented interactions at the system boundary. Give every use case
a stable ID, its own table, and at least one reciprocal requirement mapping.
Include preconditions, trigger, main flow, alternate or exception flows, and
postconditions. Maintain the canonical relationships among stories,
requirements, use cases, and tests in \ChapterRef{Traceability}.

\section[High-Level Use Case Diagram]{\hspace{0.5em}High-Level Use Case Diagram}
Show the system boundary, primary actors, and major user goals before the
detailed use-case tables. Replace the generic example with project-specific
actors and use cases while preserving the stable use-case identifiers.
Figure~\ref{Figure:UseCaseDiagram} provides the template example.

\begin{figure}[htbp]
\centering
\begin{tikzpicture}[
  actor/.style={
    draw=dsnGray,
    fill=dsnGray!8,
    rounded corners=2pt,
    minimum width=2.6cm,
    minimum height=0.8cm,
    align=center,
    font=\small
  },
  use case/.style={
    ellipse,
    draw=dsnBlue,
    fill=dsnLightBlue,
    minimum width=3.8cm,
    minimum height=1.55cm,
    inner sep=1pt,
    align=center,
    font=\small
  },
  association/.style={
    draw=dsnGray,
    line width=0.7pt
  },
  system boundary/.style={
    draw=dsnBlue,
    rounded corners=4pt,
    inner sep=0.45cm
  }
]
  \node[use case] (submit) at (0,0) {\UseCaseRef{UC-001}\\Submit a Service\\Request};
  \node[use case] (status) at (4.25,0) {\UseCaseRef{UC-002}\\View Request\\Status};
  \node[use case] (export) at (8.5,0) {\UseCaseRef{UC-003}\\Export Request\\History};

  \node[system boundary, fit=(submit)(status)(export),
    label={[font=\small\bfseries, text=dsnBlue]above:\ProjectName\ System}]
    (system) {};

  \node[actor, above=1.35cm of submit] (customer) {Customer};
  \node[actor, above=1.35cm of export] (analyst) {Operations Analyst};

  \draw[association] (customer) -- (submit);
  \draw[association] (customer) -- (status);
  \draw[association] (analyst) -- (export);
\end{tikzpicture}

\caption{Generic high-level use case diagram.}
\label{Figure:UseCaseDiagram}
\end{figure}
\index{use case diagram}

\section{Submit a Service Request}
\begin{tabularx}{\linewidth}{L{3.5cm}Y}
\toprule
\textbf{Identifier} & \UseCaseID{UC-001}\\
\textbf{Mapped requirements} & \RequirementRef{REQ-F-001}, \RequirementRef{REQ-C-001}\\
\textbf{Goal} & Submit a valid request and receive a unique identifier.\\
\textbf{Primary actor} & Customer\\
\textbf{Preconditions} & The customer can access the request form.\\
\textbf{Trigger} & The customer chooses to submit a new request.\\
\textbf{Main flow} & The customer enters required data and submits it; the system validates the data, creates the request, and displays its identifier.\\
\textbf{Alternate flow} & If required data is invalid, the system identifies the affected fields and creates no request.\\
\textbf{Postconditions} & A valid request exists with a unique identifier and an initial status.\\
\bottomrule
\end{tabularx}

\section{View Request Status}
\begin{tabularx}{\linewidth}{L{3.5cm}Y}
\toprule
\textbf{Identifier} & \UseCaseID{UC-002}\\
\textbf{Mapped requirements} & \RequirementRef{REQ-F-002}, \RequirementRef{REQ-Q-001}, \RequirementRef{REQ-D-001}, \RequirementRef{REQ-C-001}\\
\textbf{Goal} & Review the current state of a submitted request.\\
\textbf{Primary actor} & Customer\\
\textbf{Preconditions} & A request with a known identifier exists.\\
\textbf{Trigger} & The customer opens or refreshes the request status.\\
\textbf{Main flow} & The system retrieves and displays the request status and most recent update.\\
\textbf{Alternate flow} & If the identifier is unknown, the system reports that no matching request was found.\\
\textbf{Postconditions} & The request remains unchanged; its current status has been presented.\\
\bottomrule
\end{tabularx}

\section{Export Request History}
\begin{tabularx}{\linewidth}{L{3.5cm}Y}
\toprule
\textbf{Identifier} & \UseCaseID{UC-003}\\
\textbf{Mapped requirements} & \RequirementRef{REQ-B-001}, \RequirementRef{REQ-F-003}, \RequirementRef{REQ-I-001}, \RequirementRef{REQ-S-001}, \RequirementRef{REQ-D-001}, \RequirementRef{REQ-C-001}\\
\textbf{Goal} & Export request history for authorized analysis.\\
\textbf{Primary actor} & Operations analyst\\
\textbf{Preconditions} & The actor is authorized and matching requests exist.\\
\textbf{Trigger} & The actor selects an export range and format.\\
\textbf{Main flow} & The system validates authorization, generates the export, and makes the result available to the actor.\\
\textbf{Alternate flow} & If no records match, the system produces an empty result with an explanatory status.\\
\textbf{Postconditions} & An export artifact exists; source request records are unchanged.\\
\bottomrule
\end{tabularx}
